\documentclass{article}
\usepackage{amsmath,amssymb,amsthm}
\usepackage{algorithm}
\usepackage{algpseudocode}
\usepackage{fullpage}
\usepackage{hyperref}
\newtheorem{theorem}{Theorem}
\newtheorem{lemma}{Lemma}
\newtheorem{assumption}{Assumption}
\newcommand{\prox}{\operatorname{prox}}
\newcommand{\R}{\mathbb{R}}

\begin{document}

\section*{Proximal Gradient Descent (PGD)}

\section*{Mathematical Formulation}
Let $f:\R^{d}\rightarrow\R$ be a convex differentiable function with $L$-Lipschitz gradient and let $g:\R^{d}\rightarrow\R\cup\{+\infty\}$ be a proper closed convex (possibly non‑smooth) function.  
For a stepsize $\eta>0$ the \emph{proximal gradient update} is  

\[
\boxed{%
x_{k+1}= \prox_{\eta g}\!\Bigl(x_{k}-\eta \nabla f(x_{k})\Bigr)},
\qquad k=0,1,2,\dots
\]

where  

\[
\prox_{\eta g}(z)\;:=\;\arg\min_{u\in\R^{d}}\Bigl\{ g(u)+\frac{1}{2\eta}\|u-z\|^{2}\Bigr\}.
\]

\section*{Pseudocode}
\begin{algorithm}[H]
\caption{Proximal Gradient Descent}
\begin{algorithmic}[1]
\Require Convex $f$, convex $g$, stepsize $\eta\in(0,1/L]$, initial point $x_{0}\in\R^{d}$, tolerance $\varepsilon>0$, max iterations $K_{\max}$.
\For{$k=0$ \textbf{to} $K_{\max}-1$}
    \State Compute gradient $g_{k}\gets\nabla f(x_{k})$.
    \State Take proximal step 
    \[
    x_{k+1}\gets \prox_{\eta g}\bigl(x_{k}-\eta g_{k}\bigr).
    \]
    \If{$\|x_{k+1}-x_{k}\|\le \varepsilon$}
        \State \textbf{break}
    \EndIf
\EndFor
\State \textbf{return} $x_{k+1}$
\end{algorithmic}
\end{algorithm}

\section*{Dry Run Example}
Consider the scalar case $d=1$, $f(w)=(w-3)^{2}$, $g\equiv0$, $\eta=0.1$, $w_{0}=0$.

\begin{center}
\begin{tabular}{c|c|c|c}
iteration $k$ & $\nabla f(w_{k})$ & $w_{k+1}=w_{k}-\eta\nabla f(w_{k})$ & $f(w_{k+1})$\\\hline
0 & $2(w_{0}-3)=-6$ & $0-0.1(-6)=0.6$ & $(0.6-3)^{2}=5.76$\\
1 & $2(0.6-3)=-4.8$ & $0.6-0.1(-4.8)=1.08$ & $(1.08-3)^{2}=3.6864$\\
2 & $2(1.08-3)=-3.84$ & $1.08-0.1(-3.84)=1.464$ & $(1.464-3)^{2}=2.368896$\\
\end{tabular}
\end{center}
The objective value decreases monotonically, illustrating the algorithm’s behaviour.

\newpage
\section*{Assumptions}
\begin{assumption}[Convexity and Smoothness of $f$]\label{ass:convex-smooth}
$f$ is convex and differentiable with $L$-Lipschitz gradient, i.e.  
\[
\|\nabla f(x)-\nabla f(y)\|\le L\|x-y\|,\qquad\forall x,y\in\R^{d}.
\]
\end{assumption}
\begin{assumption}[Convexity of $g$]\label{ass:convex-g}
$g$ is proper, closed and convex (possibly nonsmooth). Its proximal operator $\prox_{\eta g}$ is well defined for all $\eta>0$.
\end{assumption}
\begin{assumption}[Stepsize]\label{ass:stepsize}
The stepsize satisfies $0<\eta\le 1/L$.
\end{assumption}

\section*{Main Convergence Theorem}
\begin{theorem}[Ergodic $O(1/k)$ rate]\label{thm:rate}
Let $\{x_{k}\}_{k\ge0}$ be generated by PGD under Assumptions \ref{ass:convex-smooth}–\ref{ass:stepsize} and let $x^{\star}\in\arg\min_{x}(f+g)(x)$. Then for all $K\ge1$,
\[
\frac{1}{K}\sum_{k=0}^{K-1}\bigl( f(x_{k})+g(x_{k})-f(x^{\star})-g(x^{\star})\bigr)
\;\le\;
\frac{\|x_{0}-x^{\star}\|^{2}}{2\eta K}.
\]
Consequently, there exists at least one iterate $x_{k}$, $k\in\{0,\dots,K-1\}$, satisfying  
\[
f(x_{k})+g(x_{k})-f(x^{\star})-g(x^{\star})\le \frac{\|x_{0}-x^{\star}\|^{2}}{2\eta K}.
\]
\end{theorem}

\section*{Theoretical Properties}
\begin{itemize}
\item \textbf{Stability.} The non‑expansiveness of $\prox_{\eta g}$ guarantees that the mapping $x\mapsto \prox_{\eta g}(x-\eta\nabla f(x))$ is firmly non‑expansive for $\eta\le 1/L$, preventing divergence.
\item \textbf{Hyperparameter Sensitivity.} The stepsize bound $\eta\le1/L$ is tight: larger $\eta$ may violate the descent inequality derived in the proof.
\item \textbf{Comparison to Gradient Descent.} When $g\equiv0$, $\prox_{\eta g}$ reduces to the identity map and Theorem \ref{thm:rate} collapses to the classic $O(1/k)$ rate for gradient descent on smooth convex $f$.
\item \textbf{Extension.} If $f$ is $\mu$‑strongly convex, the same proof yields a linear rate $O((1-\eta\mu)^{k})$.
\end{itemize}

\section*{Discussion}
The proof hinges on two elementary tools: (i) the descent lemma for $L$‑smooth functions and (ii) the non‑expansiveness of the proximal operator. No stochasticity is involved, so expectations do not appear. The result matches the optimal rate for first‑order methods on the class of convex composite problems.

\newpage
\section*{Preliminaries (Definitions and Lemmas)}
\begin{lemma}[Descent Lemma]\label{lem:descent}
If $f$ satisfies Assumption \ref{ass:convex-smooth}, then for any $x,y\in\R^{d}$,
\[
f(y)\le f(x)+\langle\nabla f(x),y-x\rangle+\frac{L}{2}\|y-x\|^{2}.
\]
\end{lemma}
\begin{lemma}[Non‑expansiveness of the proximal operator]\label{lem:nonexp}
For any $\eta>0$ and any $u,v\in\R^{d}$,
\[
\bigl\|\prox_{\eta g}(u)-\prox_{\eta g}(v)\bigr\|\le\|u-v\|.
\]
\end{lemma}
\begin{lemma}[Optimality condition of the proximal map]\label{lem:optcond}
For $z\in\R^{d}$ define $x^{+}= \prox_{\eta g}(z)$. Then  
\[
\frac{1}{\eta}\bigl(z-x^{+}\bigr)\in\partial g(x^{+}).
\]
\end{lemma}

\newpage
\section*{Main Proof (Step‑by‑Step Derivation)}
Let $x^{\star}\in\arg\min (f+g)$. Define the PGD update
\[
x_{k+1}= \prox_{\eta g}\!\bigl(x_{k}-\eta\nabla f(x_{k})\bigr).
\]

\noindent\textbf{[Step 1]} Apply Lemma \ref{lem:optcond} with $z_{k}=x_{k}-\eta\nabla f(x_{k})$:
\[
\frac{1}{\eta}\bigl(z_{k}-x_{k+1}\bigr)=\frac{1}{\eta}\bigl(x_{k}-\eta\nabla f(x_{k})-x_{k+1}\bigr)\in\partial g(x_{k+1}).
\]
Hence there exists $s_{k+1}\in\partial g(x_{k+1})$ such that
\[
s_{k+1}= \frac{1}{\eta}\bigl(x_{k}-x_{k+1}\bigr)-\nabla f(x_{k}). \tag{1}
\]

\noindent\textbf{[Step 2]} Convexity of $g$ gives
\[
g(x_{k+1})-g(x^{\star})\le \langle s_{k+1},\,x_{k+1}-x^{\star}\rangle . \tag{2}
\]

\noindent\textbf{[Step 3]} Plug the expression of $s_{k+1}$ from (1) into (2):
\[
g(x_{k+1})-g(x^{\star})\le
\Big\langle\frac{1}{\eta}(x_{k}-x_{k+1})-\nabla f(x_{k}),\,x_{k+1}-x^{\star}\Big\rangle .
\]

\noindent\textbf{[Step 4]} Rearrange the inner product:
\[
\begin{aligned}
g(x_{k+1})-g(x^{\star})
&\le\frac{1}{\eta}\langle x_{k}-x_{k+1},x_{k+1}-x^{\star}\rangle
-\langle\nabla f(x_{k}),x_{k+1}-x^{\star}\rangle\\
&=\frac{1}{2\eta}\Bigl(\|x_{k}-x^{\star}\|^{2}-\|x_{k+1}-x^{\star}\|^{2}
-\|x_{k+1}-x_{k}\|^{2}\Bigr)
-\langle\nabla f(x_{k}),x_{k+1}-x^{\star}\rangle .
\end{aligned}
\tag{3}
\]

\noindent\textbf{[Step 5]} Apply the descent lemma (Lemma \ref{lem:descent}) with $x=x_{k}$ and $y=x_{k+1}$:
\[
f(x_{k+1})\le f(x_{k})+\langle\nabla f(x_{k}),x_{k+1}-x_{k}\rangle+\frac{L}{2}\|x_{k+1}-x_{k}\|^{2}.
\tag{4}
\]

\noindent\textbf{[Step 6]} Add inequality (3) and (4). The terms $\langle\nabla f(x_{k}),x_{k+1}-x_{k}\rangle$ cancel partially:
\[
\begin{aligned}
&\bigl[f(x_{k+1})+g(x_{k+1})\bigr]-\bigl[f(x^{\star})+g(x^{\star})\bigr] \\
&\le f(x_{k})-f(x^{\star})+\frac{L}{2}\|x_{k+1}-x_{k}\|^{2}
-\langle\nabla f(x_{k}),x_{k+1}-x^{\star}\rangle\\
&\quad+\frac{1}{2\eta}\Bigl(\|x_{k}-x^{\star}\|^{2}-\|x_{k+1}-x^{\star}\|^{2}
-\|x_{k+1}-x_{k}\|^{2}\Bigr).
\end{aligned}
\tag{5}
\]

\noindent\textbf{[Step 7]} Observe that 
\[
-\langle\nabla f(x_{k}),x_{k+1}-x^{\star}\rangle
= -\langle\nabla f(x_{k}),x_{k+1}-x_{k}\rangle-\langle\nabla f(x_{k}),x_{k}-x^{\star}\rangle .
\]
Insert this into (5) and use (4) again to replace $-\langle\nabla f(x_{k}),x_{k+1}-x_{k}\rangle$ by 
$f(x_{k})-f(x_{k+1})-\frac{L}{2}\|x_{k+1}-x_{k}\|^{2}$:
\[
\begin{aligned}
&\bigl[f(x_{k+1})+g(x_{k+1})\bigr]-\bigl[f(x^{\star})+g(x^{\star})\bigr] \\
&\le \frac{1}{2\eta}\bigl(\|x_{k}-x^{\star}\|^{2}-\|x_{k+1}-x^{\star}\|^{2}\bigr)
-\Bigl(\frac{1}{2\eta}-\frac{L}{2}\Bigr)\|x_{k+1}-x_{k}\|^{2}.
\end{aligned}
\tag{6}
\]

\noindent\textbf{[Step 8]} By the stepsize condition $\eta\le 1/L$, the coefficient of $\|x_{k+1}-x_{k}\|^{2}$ in (6) is non‑negative; dropping this non‑negative term yields the fundamental descent inequality
\[
\boxed{
f(x_{k+1})+g(x_{k+1})-f(x^{\star})-g(x^{\star})
\le \frac{1}{2\eta}\Bigl(\|x_{k}-x^{\star}\|^{2}-\|x_{k+1}-x^{\star}\|^{2}\Bigr).
}
\tag{7}
\]

\noindent\textbf{[Step 9]} Sum (7) from $k=0$ to $K-1$ telescopically:
\[
\sum_{k=0}^{K-1}\bigl[f(x_{k+1})+g(x_{k+1})-f(x^{\star})-g(x^{\star})\bigr]
\le \frac{1}{2\eta}\bigl(\|x_{0}-x^{\star}\|^{2}-\|x_{K}-x^{\star}\|^{2}\bigr)
\le \frac{\|x_{0}-x^{\star}\|^{2}}{2\eta}.
\tag{8}
\]

\noindent\textbf{[Step 10]} Divide both sides of (8) by $K$ and note that the left‑hand side is an average over the first $K$ iterates (re‑indexing $k+1\to k$) to obtain the bound stated in Theorem \ref{thm:rate}.

\section*{Rate Derivation (With Constants)}
From (7) we have for every $k$,
\[
f(x_{k})+g(x_{k})-f(x^{\star})-g(x^{\star})
\le \frac{\|x_{k-1}-x^{\star}\|^{2}}{2\eta}
-\frac{\|x_{k}-x^{\star}\|^{2}}{2\eta}.
\]
Summing and using $\|x_{K}-x^{\star}\|^{2}\ge0$ leads to the explicit constant
\[
\frac{1}{K}\sum_{k=0}^{K-1}\bigl(f(x_{k})+g(x_{k})-f(x^{\star})-g(x^{\star})\bigr)
\le \frac{\|x_{0}-x^{\star}\|^{2}}{2\eta K}.
\]
Thus the convergence rate is $O(1/K)$ with constant $C=\|x_{0}-x^{\star}\|^{2}/(2\eta)$.

\section*{Special Cases}
\begin{itemize}
\item \textbf{Pure Gradient Descent ($g\equiv0$).} $\prox_{\eta g}$ becomes the identity, and inequality (7) reduces to the classic smooth convex bound
\[
f(x_{k+1})-f(x^{\star})\le\frac{1}{2\eta}\bigl(\|x_{k}-x^{\star}\|^{2}-\|x_{k+1}-x^{\star}\|^{2}\bigr),
\]
which yields the same $O(1/K)$ rate.
\item \textbf{Strongly Convex $f+g$.} If $f+g$ is $\mu$‑strongly convex, then from (7) one obtains
\[
\|x_{k+1}-x^{\star}\|^{2}\le\bigl(1-\eta\mu\bigr)\|x_{k}-x^{\star}\|^{2},
\]
implying a linear convergence $O\bigl((1-\eta\mu)^{k}\bigr)$.
\item \textbf{Indicator $g=\iota_{C}$ (projection onto a convex set $C$).} The proximal step reduces to Euclidean projection $P_{C}(\cdot)$. The algorithm becomes the projected gradient method, and the same $O(1/K)$ guarantee holds.
\end{itemize}

\end{document}